\chapter{Delta Lake Optimization: OPTIMIZE, Z-ORDER, and Liquid Clustering}
\index{Delta Lake!optimization}\index{OPTIMIZE}\index{Z-ORDER}\index{liquid clustering}%
\index{small-file problem}\index{data skipping}\index{deletion vectors}\index{VACUUM}\index{predictive optimization}%
\begin{objectives}
\begin{itemize}
    \item Explain how file-level statistics enable data skipping, and why the small-file problem defeats it.
    \item Compact tables with \texttt{OPTIMIZE} and cluster them with \texttt{ZORDER BY}.
    \item Choose liquid clustering (\texttt{CLUSTER BY}) for new tables and articulate when partitioning still applies.
    \item Use deletion vectors to make row-level deletes cheap, and auto-optimize versus scheduled compaction deliberately.
    \item Operate \texttt{VACUUM} and predictive optimization within a retention policy.
    \item Benchmark layout choices with controlled before/after measurements.
\end{itemize}
\end{objectives}

\begin{crosscloud}
Physical layout optimization exists on every analytical platform; the knobs differ, the physics
does not.

\vspace{2mm}
{\footnotesize
\begin{tabular}{@{}p{3.0cm}p{3.4cm}p{3.2cm}p{3.2cm}@{}}
\toprule
\textbf{Capability} & \textbf{Databricks (Delta)} & \textbf{AWS} & \textbf{Azure / GCP / Fabric} \\
\midrule
Compaction & \texttt{OPTIMIZE} & Athena/Iceberg \texttt{OPTIMIZE} & \texttt{OPTIMIZE} (Fabric) / BigQuery managed \\
Multi-dim clustering & \texttt{ZORDER BY} & Redshift sort keys & Synapse columnstore / Fabric V-Order+ZORDER \\
Rewritable-free layout & Liquid clustering & Iceberg sort order evolution & BigQuery clustering / clustering keys (Snowflake) \\
Row-level deletes & Deletion vectors & Iceberg position deletes & Native / deletion vectors (Fabric) \\
Auto-maintenance & Predictive optimization & Redshift Auto WLM / S3 Tables & Automatic / BigQuery autoclustering \\
Pruning metadata & File stats, bloom filters & Zone maps (Redshift) & Columnstore segments / block metadata \\
\bottomrule
\end{tabular}}

\vspace{1mm}
\textbf{MongoDB} answers the same question with indexes and the working set rather than file
layout---an operational-store solution to an analytical problem it was not designed for.
\end{crosscloud}

\section{Week 5 Roadmap and Mental Model}
Chapter 2 made tables correct; Chapter 5 makes them fast. Analytical query speed on a lakehouse is dominated by one question: \emph{how many bytes must be read to answer this query?} Everything in this chapter---compaction, clustering, statistics, deletion vectors---exists to shrink that number \cite{databricksDeltaOptimize}.

The mental model: \textbf{a Delta scan is a funnel}. The query's predicates are evaluated first against the transaction log's per-file statistics (min/max per column, null counts), discarding files that cannot match; surviving files are opened, and column pruning plus row-group statistics discard more; only then are rows decoded. The optimizer's job is to make the funnel narrow early. The engineer's job is to arrange data so that statistics \emph{can} discard---which is what clustering does.

\begin{definitionbox}
\textbf{Data skipping} is the elimination of data files (and row groups) from a scan based on metadata alone, without reading row content. It works only when the data layout correlates file contents with the query's filter columns---the property that Z-ORDER and liquid clustering create.
\end{definitionbox}

\begin{keyterms}
\textbf{Small-file problem}: thousands of tiny files whose metadata overhead and per-file cost dominate query time. \textbf{Z-ORDER}: multi-dimensional clustering that co-locates rows similar on several columns in the same files. \textbf{Liquid clustering}: Delta's self-managing layout where clustering keys are metadata, not directory structure. \textbf{Deletion vector}: a bitmap marking deleted rows so deletes need not rewrite files. \textbf{Predictive optimization}: Databricks-managed automatic \texttt{OPTIMIZE}/\texttt{ANALYZE}/\texttt{VACUUM} for Unity Catalog tables.
\end{keyterms}

\section{Monday: File Layout, Statistics, and the Small-File Problem}
\subsection{How a Delta Scan Reads}
Each commit records added files with per-column min/max statistics (for the first 32 columns by default). A query like \texttt{WHERE order\_date = '2026-08-01'} consults the log and opens only files whose statistics admit that date. Bloom filters (opt-in) add point-lookup skipping for high-cardinality columns such as \texttt{order\_id}.

\subsection{The Small-File Problem}
Every file costs: a listing/metadata entry, a footer read, a task slot, and open/close overhead. At thousands of 1 MB files, metadata work dwarfs data work; queries crawl, and \texttt{OPTIMIZE} itself becomes slow. Small files come from streaming micro-batches, over-partitioned writes, and un-coalesced outputs (Chapter 4). The fix is compaction; the prevention is write cadence \cite{databricksDeltaOptimize}.

\begin{pitfall}
\textbf{Optimizing while streaming into the same table.} A compaction job rewriting files that a stream is actively appending creates commit conflicts and wasted work. Coordinate schedules---triggered streams followed by \texttt{OPTIMIZE}---or use predictive optimization, which is conflict-aware.
\end{pitfall}

\section{Tuesday: OPTIMIZE and Z-ORDER}
\subsection{Compaction}
\texttt{OPTIMIZE table} rewrites small files into $\sim$1 GB targets. It is incremental (only unoptimized files are touched) and recorded in history like any commit. Run it on a schedule matched to write cadence: daily for batch-loaded tables, after each trigger group for streaming sinks.

\subsection{Z-ORDER}
\texttt{ZORDER BY (col1, col2)} adds multi-dimensional clustering: the compaction sorts data so rows similar on the given columns co-locate, making file-level statistics sharply effective for filters on those columns.

\begin{lstlisting}[language=SQL,caption={Compaction and Z-ORDER, with evidence collection.},label={lst:ch5-optimize}]
-- Baseline first: file count and a representative selective query
SELECT region, COUNT(*) c, SUM(net_amount) AS rev
FROM de_training.silver.events
WHERE event_date = DATE '2026-08-01' AND customer_id = 4217
GROUP BY region;

OPTIMIZE de_training.silver.events;                          -- compaction only
OPTIMIZE de_training.silver.events
  ZORDER BY (customer_id, event_date);                       -- + clustering

DESCRIBE DETAIL de_training.silver.events;                   -- numFiles, sizeInBytes
DESCRIBE HISTORY de_training.silver.events;                  -- OPTIMIZE commits
\end{lstlisting}

\begin{tipbox}
Choose Z-ORDER columns from \emph{actual query patterns}: high-cardinality columns used in selective equality or range filters. Z-ORDER on low-cardinality columns (like \texttt{region} with 5 values) wastes the sort; the table's partition columns rarely need to be Z-ORDER columns too.
\end{tipbox}

\section{Wednesday: Liquid Clustering, Deletion Vectors, VACUUM}
\subsection{Liquid Clustering}
Liquid clustering (\texttt{CLUSTER BY}) makes clustering a table \emph{property} rather than a physical directory layout. Unlike Hive-style partitioning, clustering keys can change (\texttt{ALTER TABLE ... CLUSTER BY (...)}) without rewriting the table; unlike Z-ORDER, clustering is maintained incrementally on write and optimize, and it composes with automatic file sizing \cite{databricksLiquidClustering}.

\begin{lstlisting}[language=SQL,caption={Liquid clustering: creation, change, and the partition migration.},label={lst:ch5-liquid}]
-- New table
CREATE TABLE de_training.gold.fact_orders_lc (
  order_id BIGINT, customer_id BIGINT, order_date DATE,
  region STRING, net_amount DECIMAL(12,2))
CLUSTER BY (order_date, customer_id);

-- Existing table: opt in without a rewrite
ALTER TABLE de_training.gold.fact_orders
  CLUSTER BY (order_date, customer_id);

-- Change clustering keys later: metadata operation
ALTER TABLE de_training.gold.fact_orders
  CLUSTER BY (region, order_date);
\end{lstlisting}

\begin{definitionbox}
\textbf{Partitioning} encodes a column into directory paths---fast pruning, but immutable once written and dangerous at high cardinality. \textbf{Liquid clustering} records clustering in the log---similar pruning benefits, re-definable keys, no small-partition hazard. The default for new large Delta tables is liquid clustering; partitioning survives for genuinely coarse, stable dimensions like \texttt{country} or \texttt{business\_date} at low cardinality.
\end{definitionbox}

\subsection{Deletion Vectors}
With deletion vectors enabled, \texttt{DELETE}/\texttt{UPDATE}/\texttt{MERGE} mark rows deleted via a bitmap stored alongside the file instead of rewriting it \cite{databricksDeletionVectors}. Reads apply the bitmap on the fly; a later \texttt{OPTIMIZE} physically purges. Enable deletion vectors on tables with frequent row-level mutations (CDC-heavy Silvers, GDPR deletes).

\subsection{VACUUM and Predictive Optimization}
\texttt{VACUUM table RETAIN 168 HOURS} removes unreferenced files past the horizon (Chapter 2 covered the recovery implications). With Unity Catalog managed tables, \textbf{predictive optimization} runs \texttt{OPTIMIZE}, \texttt{ANALYZE}, and \texttt{VACUUM} automatically where Databricks' telemetry shows benefit---on serverless, at no separate job cost. Enable it and spend your scheduling effort on the exceptions.

\section{Thursday Hands-on Project: The Small-File Benchmark}
Build a 50M-row events table degraded by tiny files, then measure the effect of compaction, Z-ORDER, and liquid clustering on the same selective query.
\index{hands-on lab}\index{OPTIMIZE!lab}\index{liquid clustering!lab}

\begin{lstlisting}[language=Python,caption={Week 5 lab: degrade, then optimize, with controlled measurement.},label={lst:ch5-lab}]
from pyspark.sql import functions as F
import time

# 1. Generate 50M rows as ~2000 tiny files (the streaming-aftermath simulation)
base = (spark.range(0, 50_000_000)
        .withColumn("event_date", F.date_add(F.lit("2026-01-01"),
                                             (F.rand(7) * 240).cast("int")))
        .withColumn("customer_id", (F.rand(7) * 500_000).cast("long"))
        .withColumn("amount", F.rand(7) * 100))
base.repartition(2000).write.mode("overwrite") \
    .saveAsTable("de_training.silver.events_raw")

def bench(table, where, label):
    spark.sql(f"SELECT * FROM {table} LIMIT 1").count()  # warm metadata
    t0 = time.time()
    spark.sql(f"SELECT COUNT(*) c, SUM(amount) s FROM {table} WHERE {where}") \
         .collect()
    return (label, round(time.time() - t0, 2))

sel = "event_date = '2026-05-01' AND customer_id = 4217"
results = [bench("de_training.silver.events_raw", sel, "baseline_small_files")]

# 2. Compaction + Z-ORDER copy
spark.sql("""CREATE TABLE de_training.silver.events_zo
             AS SELECT * FROM de_training.silver.events_raw""")
spark.sql("""OPTIMIZE de_training.silver.events_zo
             ZORDER BY (customer_id, event_date)""")
results.append(bench("de_training.silver.events_zo", sel, "optimize_zorder"))

# 3. Liquid-clustered copy
spark.sql("""CREATE TABLE de_training.silver.events_lc
             CLUSTER BY (event_date, customer_id)
             AS SELECT * FROM de_training.silver.events_raw""")
spark.sql("OPTIMIZE de_training.silver.events_lc")
results.append(bench("de_training.silver.events_lc", sel, "liquid_clustered"))

spark.createDataFrame(results, ["variant", "seconds"]) \
     .write.mode("overwrite").saveAsTable("de_training.silver.week05_results")
\end{lstlisting}

\subsection*{Deliverables}
\begin{itemize}
    \item The three variants (raw, Z-ORDER, liquid) with \texttt{DESCRIBE DETAIL} file counts.
    \item The \texttt{week05\_results} timing table with cold-run and warm-run columns.
    \item A one-paragraph interpretation: which layout won, and why the statistics made it so.
\end{itemize}

\subsection*{Acceptance Criteria}
\begin{itemize}
    \item Identical result values across all three variants (layout changes must not change data).
    \item File counts documented before and after compaction.
    \item Benchmarks repeated three times with the median reported (first-run cold effects excluded).
    \item \texttt{VACUUM} not run below 168-hour retention at any point.
\end{itemize}

\subsection*{Stretch Goals}
\begin{itemize}
    \item Add a bloom filter on \texttt{customer\_id} and measure the point-query delta.
    \item Enable deletion vectors, delete 1\% of rows, and compare delete cost with/without.
    \item Repeat the benchmark on a partitioned variant and explain where partitioning beat clustering.
\end{itemize}

\section{Friday Use Case: Remediating Two Million Tiny Files}
A nightly streaming sink has appended micro-batches to \texttt{silver.events} for six months. Analysts report multi-minute scans; the metadata listing alone takes minutes; \texttt{OPTIMIZE} times out. The data is correct---the layout has collapsed.
\index{small-file problem!remediation}\index{streaming!compaction}

\subsection{Symptoms and Diagnosis}
Evidence collected: \texttt{DESCRIBE DETAIL} shows 2.1M files averaging 0.8 MB; stage summaries show thousands of one-file tasks; the driver spends most of each query planning against an enormous file list. The streaming trigger fires every 30 seconds; no compaction has ever run.

\subsection{Remediation Strategy}
Two phases: \textbf{repair}, then \textbf{prevent}.

\begin{figure}[H]
    \centering
    \resizebox{\textwidth}{!}{%
    \begin{tikzpicture}[
        tiny/.style={rectangle, draw=red!60!black, thick, fill=red!8, align=center, minimum width=0.85cm, minimum height=0.4cm, font=\tiny},
        big/.style={rectangle, rounded corners, draw=codegreen!60!black, thick, fill=green!8, align=center, minimum width=2.4cm, minimum height=0.85cm, font=\small\bfseries},
        proc/.style={rectangle, rounded corners, draw=primary, thick, fill=primary!7, align=center, minimum width=3.0cm, minimum height=0.95cm, font=\small\bfseries},
        flow/.style={-{Latex[length=2mm]}, draw=primary, thick}
    ]
        \node[tiny] (f1) at (0,0.6) {}; \node[tiny] (f2) at (0.9,0.6) {};
        \node[tiny] (f3) at (0,0.0) {}; \node[tiny] (f4) at (0.9,0.0) {};
        \node[tiny] (f5) at (0,-0.6) {}; \node[tiny] (f6) at (0.9,-0.6) {};
        \node[font=\scriptsize\itshape] at (0.45,-1.2) {2.1M tiny files};
        \node[proc] (pause) at (3.6,0) {Pause stream\\snapshot history};
        \node[proc] (compact) at (7.4,0) {Deep \texttt{OPTIMIZE}\\(batched by date)};
        \node[big] (healthy) at (11.0,0) {Healthy files\\$\sim$1 GB each};
        \node[proc] (prevent) at (14.6,0) {Liquid clustering\\+ trigger cadence\\+ predictive opt.};
        \draw[flow] (1.4,0) -- (pause);
        \draw[flow] (pause) -- (compact);
        \draw[flow] (compact) -- (healthy);
        \draw[flow] (healthy) -- (prevent);
    \end{tikzpicture}%
    }
    \caption{Small-file remediation: pause, snapshot, batch compaction, then prevent recurrence with clustering and write-cadence changes.}
    \label{fig:ch5-remediation}
\end{figure}

\textbf{Repair.} Pause or throttle the stream (triggered mode with a longer interval), snapshot the history, and run compaction in date-banded batches so each \texttt{OPTIMIZE} operates on a tractable file set. Verify counts against the pre-repair baseline; only then resume.

\textbf{Prevent.} Convert the table to liquid clustering on the dominant filter columns; move the stream from 30-second triggers to 5-minute micro-batches (or \texttt{availableNow} batches); enable predictive optimization; add a weekly file-count guardrail alert.

\begin{pitfall}
\textbf{VACUUM in the middle of repair.} During remediation you are one mistake away from needing history. Vacuum only after the repaired table has been validated and the stream has run cleanly against it for a full retention window.
\end{pitfall}

\subsection{Risks and Mitigations}
\begin{table}[H]
\centering
\small
\begin{tabular}{@{}p{4.6cm}p{7.6cm}@{}}
\toprule
\textbf{Risk} & \textbf{Mitigation} \\
\midrule
Compaction job becomes the new bottleneck & Date-banded runs; predictive optimization absorbs routine work \\
Layout change regresses unmeasured queries & Benchmark harness covers the top five query patterns \\
Small files return via a new producer & File-count guardrail alert on every hot table \\
VACUUM mid-incident destroys evidence & Freeze VACUUM during open incidents; runbook says so \\
Clustering keys chosen by intuition & Keys derived from query-history analysis, reviewed quarterly \\
\bottomrule
\end{tabular}
\caption{Week 5 scenario risks and mitigations.}
\label{tab:ch5-risks}
\end{table}

\section{Design Decision Guide}
\index{decision guide!Week 5}%
Layout decisions age with the table; choose for the query patterns you can verify, not the ones you imagine.

\begin{table}[H]
\centering
\small
\begin{tabular}{@{}p{3.4cm}p{4.4cm}p{4.6cm}@{}}
\toprule
\textbf{Decision} & \textbf{Choose the first when} & \textbf{Choose the second when} \\
\midrule
Liquid clustering vs.\ partition+Z-ORDER & New large tables; evolving filter columns & Legacy tables you cannot migrate yet \\
Predictive optimization vs.\ scheduled OPTIMIZE & UC managed tables on serverless & External tables or custom cadence needs \\
Deletion vectors on vs.\ off & Row-level mutation is frequent & Append-mostly tables with scan-sensitive SLAs \\
Bloom filter vs.\ statistics only & High-cardinality point lookups & Range-dominated queries \\
Daily vs.\ post-trigger compaction & Batch-loaded tables & Streaming sinks with micro-batches \\
Partition vs.\ cluster on date & Coarse day-grain filters and lifecycle drops & Selective queries on mixed columns \\
\bottomrule
\end{tabular}
\caption{Week 5 decision guide.}
\label{tab:ch5-decisions}
\end{table}

Benchmark before believing: the same clustering choice can win on one table and waste money on another. Chapter 5's lab harness is the template for every layout decision you will ever make.

\section{Operational Guidance for Week 5}
\index{operational guidance!Week 5}%
\subsection{Performance and Reliability}
Optimization is a cadence, not an event: set \texttt{OPTIMIZE} frequency from write volume (daily for batch-loaded tables, after each trigger window for streaming sinks) and review file counts monthly against a growth guardrail. Prefer predictive optimization on Unity Catalog managed tables so maintenance follows measured benefit; keep scheduled jobs for external tables and exceptions. Track the benchmark trio---file count, selective-query latency, scan bytes---before and after every layout change.

\subsection{Security and Governance Checklist}
\begin{itemize}
    \item \texttt{ALTER TABLE ... CLUSTER BY} and \texttt{OPTIMIZE} rights restricted to the platform group.
    \item Layout changes (clustering keys, partition retirement) documented in the change log with benchmark evidence.
    \item Deletion vectors enabled on tables with GDPR/row-delete obligations---and the purge cadence documented.
    \item \texttt{VACUUM} schedules reviewed against legal hold and recovery requirements quarterly.
\end{itemize}

\subsection{Troubleshooting Playbook}
\begin{table}[H]
\centering
\small
\begin{tabular}{@{}p{3.6cm}p{4.2cm}p{4.8cm}@{}}
\toprule
\textbf{Symptom} & \textbf{Likely cause} & \textbf{First action} \\
\midrule
OPTIMIZE conflicts with writers & Compaction racing the stream & Stagger schedules or enable predictive optimization \\
Z-ORDER shows no gain & Low-cardinality or rarely-filtered columns chosen & Re-derive columns from actual query predicates \\
File count not dropping & Retention keeping old files / optimize not scheduled & Check history for OPTIMIZE commits; review VACUUM \\
Reads slower after enabling DVs & Deletion bitmap merge overhead & Run OPTIMIZE to physically purge; review delete rate \\
\bottomrule
\end{tabular}
\caption{Week 5 troubleshooting quick reference.}
\label{tab:ch5-trouble}
\end{table}

\section{Interview Q\&A}
\subsection{Concepts and Trade-offs}
\begin{enumerate}
    \item \textbf{Q: Why do small files hurt even when total bytes are the same?}\\
    A: Per-file costs---listing, footer reads, task scheduling, open/close---scale with file count, not bytes. At millions of files, metadata work dominates and the driver becomes the bottleneck.
    \item \textbf{Q: Z-ORDER versus liquid clustering?}\\
    A: Both co-locate similar rows for data skipping. Liquid clustering keeps keys as metadata (changeable without rewrite, maintained incrementally, no partition-column commitment); Z-ORDER applies at compaction time on a fixed layout. New tables: liquid; legacy: Z-ORDER until migrated.
    \item \textbf{Q: When does classic partitioning still win?}\\
    A: Coarse, stable, low-cardinality dimensions that almost every query filters (e.g., \texttt{country}, business date at day grain), where directory pruning is exact and lifecycle operations drop whole directories.
    \item \textbf{Q: What do deletion vectors change about a DELETE?}\\
    A: The delete writes a bitmap instead of rewriting files---seconds instead of minutes---but every reader pays a small merge cost until the next OPTIMIZE purges physically.
    \item \textbf{Q: Your OPTIMIZE job itself is slow. Why?}\\
    A: Usually because it is fighting the same small-file mountain (schedule it more often, in bands) or competing with active writers (coordinate schedules or use predictive optimization).
    \item \textbf{Q: Does \texttt{OPTIMIZE} change data or only layout?}\\
    A: Layout only---it is a commit like any other, visible in history, and your reconciliation counts must be identical before and after. A row-count change after OPTIMIZE is a bug report, not a curiosity.
    \item \textbf{Q: How would you convince a skeptic that liquid clustering helped?}\\
    A: The lab's evidence table: same query, same data, controlled medians across cold and warm runs, plus the file-level statistics story. Feelings are not benchmarks.
    \item \textbf{Q: Does liquid clustering remove the need to think about layout?}\\
    A: It removes the \emph{irreversible} decision, not the decision: clustering keys still determine skipping quality, and they should still come from measured query patterns.
\end{enumerate}

\section{Further Reading}
The Delta optimization guide covers compaction, Z-ORDER, bloom filters, and file sizing \cite{databricksDeltaOptimize}; the liquid clustering documentation explains the layout model and migration \cite{databricksLiquidClustering}; deletion vectors are covered in \cite{databricksDeletionVectors}. For the storage fundamentals underneath, revisit the Delta paper \cite{armbrust2020delta} and the Parquet format docs \cite{apacheParquet}. Spark-level tuning complements this chapter \cite{sparkSQLTuning}.

\section*{Key Takeaways}
\begin{itemize}
    \item Query speed on a lakehouse is a bytes-scanned problem; statistics-based skipping is the mechanism.
    \item The small-file problem is a write-cadence disease; compaction is the treatment, clustering the prevention.
    \item Liquid clustering is the default layout decision for new large tables; partitioning survives for coarse stable dimensions.
    \item Deletion vectors make row-level mutation cheap at the cost of reader-side merge until compaction.
    \item Benchmark layout changes with controlled medians, and never vacuum below your recovery horizon.
\end{itemize}

\section{Exercises}
\begin{enumerate}
    \item \textbf{(Conceptual)} Explain data skipping to a PM in three sentences without using the word ``statistics.''
    \item \textbf{(Conceptual)} Why is Z-ORDER on a 5-value column wasted work?
    \item \textbf{(Hands-on)} Complete the Thursday lab; capture the results table and one \texttt{DESCRIBE DETAIL} per variant.
    \item \textbf{(Hands-on)} Write the date-banded compaction script for the Friday scenario, with per-band logging.
    \item \textbf{(Hands-on)} Enable predictive optimization on one catalog and find the evidence of its runs in system tables.
    \item \textbf{(Design)} Choose partition/cluster/Z-ORDER strategies for: (a) clickstream by day, (b) IoT by device, (c) orders by country and date. Justify each.
    \item \textbf{(Operations)} Write the alert rule and runbook for ``file count grows 20\% week-over-week.''
    \item \textbf{(Architecture)} When would you disable deletion vectors on a table? Give a workload where the reader-side merge cost matters.
\end{enumerate}
